\medskip

\begin{minipage}{0.48\linewidth}
On considère le cube ABCDEFGH d'arête 1.

On appelle I le point d'intersection du plan (GBD) avec la droite (EC).

L'espace est rapporté au repère orthonormé $\left(\text{A}~;~\vect{\text{AB}},~\vect{\text{AD}},~\vect{\text{AE}}\right)$.
\end{minipage} \hfill
\begin{minipage}{0.48\linewidth}
\psset{unit=1cm}
\begin{pspicture}(7.3,5.7)
\psframe(0.8,0.3)(4.8,4.3)
\psline(4.8,0.3)(6.3,1.3)(6.3,5.3)(4.8,4.3)
\psline(6.3,5.3)(2.3,5.3)(0.8,4.3)
\psline(4.8,0.3)(6.3,5.3)
\psline[linestyle=dashed](0.8,0.3)(2.3,1.3)(6.3,1.3)
\psline[linestyle=dashed](2.3,1.3)(2.3,5.3)
\psline[linestyle=dotted,linewidth=1.2pt](4.8,0.3)(2.3,1.3)(6.3,5.3)
\psline[linestyle=dotted,linewidth=1.2pt](0.8,4.3)(6.3,1.3)
\uput[d](0.8,0.3){A} \uput[d](4.8,0.3){B} \uput[dr](6.3,1.3){C} \uput[d](2.3,1.3){D}
\uput[u](0.8,4.3){E} \uput[u](4.8,4.3){F} \uput[u](6.3,5.3){G} \uput[u](2.3,5.3){H}
\uput[dl](4.5,2.27){I} \uput[dl](3.55,0.8){J}
\psdots(4.5,2.27)(3.55,0.8)
\end{pspicture}
\end{minipage}

\medskip

\begin{enumerate}
\item 
On a:
E\;$\begin{pmatrix} 0\\0\\1\end{pmatrix}$,
C\;$\begin{pmatrix} 1\\1\\0\end{pmatrix}$ et
G\;$\begin{pmatrix} 1\\1\\1\end{pmatrix}$.

\item La droite (EC) a pour vecteur directeur 
$\vectt{EC}\;\begin{pmatrix} 1-0\\1-0\\0-1\end{pmatrix}$
soit $\begin{pmatrix} 1\\1\\-1\end{pmatrix}$.

Donc la droite (EC) a pour représentation paramétrique;

$\left\{
\begin{array}{l !{=} l l}
x & x_{\text E} + 1\times t\\
y & y_{\text E} + 1\times t & t\in\R\\
z & z_{\text E} + (-1)\times t
\end{array}
\right.$
soit
$\left\{
\begin{array}{l !{=} r l}
x & t\\
y & t & t\in\R\\
z & 1- t
\end{array}
\right.$

\item 
Les points G, B et D ne sont pas alignés donc (GDB) est un plan de vecteurs directeurs $\vectt{GB}$ et $\vectt{BD}$.

$\vectt{GB}$ a pour coordonnées
$\begin{pmatrix} 1-1\\0-1\\0-1 \end{pmatrix}
=
\begin{pmatrix}0\\-1\\-1 \end{pmatrix}$
et
$\vectt{BD}$ a pour coordonnées
$\begin{pmatrix} 0-1\\1-0\\0-0 \end{pmatrix}
=
\begin{pmatrix}-1\\1\\0 \end{pmatrix}$.

$\vectt{EC}.\vectt{GB}=1\times 0 + 1\times (-1) + (-1)\times (-1)=0-1+1=0$ donc $\vectt{EC}\perp\vectt{GB}$

$\vectt{EC}.\vectt{BD}=1\times (-1) + 1\times 1 + (-1)\times 0=-1+1+0=0$ donc $\vectt{EC}\perp\vectt{BD}$

Le vecteur $\vectt{EC}$ est orthogonal à deux vecteurs directeurs du plan (GDB) donc la droite (EC) est orthogonale au plan (GDB).

\item 
	\begin{enumerate}
		\item 
La droite (EC) est orthogonale au plan (GDB) donc tout vecteur directeur de la droite (EC) est un vecteur normal au plan (GDB), en particulier le vecteur $\vectt{EC}$ de coordonnées $(1\;;\;1\;;\;-1)$.
Donc le plan (GDB) a une équation de la forme $x+y-z+d=0$ où $d$ est un réel à déterminer.

Le point B appartient au plan (GDB) donc $x_{\text B} + y_{\text B} -z_{\text B}+d=0$; autrement dit: $1+0-0+d=0$ donc $d=-1$.

Le plan (GDB) a donc pour équation: $x+y-z-1=0$.
			
		\item 
On appelle I le point d'intersection du plan (GBD) avec la droite (EC).

Donc les coordonnées de I sont solutions du système:
$\left \{
\begin{array}{r !{=} l}
x	& t\\
y & t\\
z & 1-t\\
x+y-z-1 & 0
\end{array}
\right .$

On a donc: $t+t-(1-t)-1=0$ donc $t=\dfrac{2}{3}$.

On en déduit les coordonnées de I:
$x_{\text I}=t=\dfrac{2}{3}$, $y_{\text I}=t=\dfrac{2}{3}$ et $z_{\text I}=1-t=\dfrac{1}{3}$,
soit $\left( \dfrac{2}{3}\;;\;\dfrac{2}{3}\;;\;\dfrac{1}{3}\right )$. 
		
		\item 
La droite (EC) est orthogonale au plan (GDB) donc la distance du point E au plan (GDB) est la longueur EI.

$\text{EI}^2 = \left ( \dfrac{2}{3} - 0\right )^2 + \left ( \dfrac{2}{3} - 0\right )^2 + \left ( \dfrac{1}{3} - 1\right )^2 = \dfrac{4}{9}+\dfrac{4}{9}+\dfrac{4}{9}=\dfrac{12}{9}$;
donc $\text{EI} = \ds\sqrt{\dfrac{12}{9}} = \dfrac{\ds\sqrt{12}}{3}=\dfrac{2\sqrt{3}}{3}$

La distance du point I au plan (GDB) est donc de $\dfrac{2\sqrt{3}}{3}$.
	\end{enumerate}
\item 
	\begin{enumerate}
		\item
Les côtés [BD], [DG] et [GB] du triangle BDG sont les diagonales respectives des carrés ABCD, CDHG et BCGF qui sont superposables; donc leurs diagonales ont la même longueur. Ce qui veut dire que le triangle BDG est équilatéral.

		\item 
		On appelle J le milieu du segment [BD]; donc le segment [GJ] est la médiane issue de G dans le triangle BDG. Or ce triangle est équilatéral donc [GJ] est la hauteur issue de G dans ce triangle.
L'aire du triangle BDG est donc:
$\mathcal A = \dfrac{\text{BD}\times \text{GJ}}{2}$.

B\;$\begin{pmatrix} 1\\ 0\\ 0\\\end{pmatrix}$
et
D\;$\begin{pmatrix} 0\\ 1\\ 0\\\end{pmatrix}$
donc
J\;$\begin{pmatrix} \frac{1}{2}\\ \frac{1}{2}\\ 0\\\end{pmatrix}$

$\text{GJ}^2 = \left ( \dfrac{1}{2} -1\right )^2 + \left ( \dfrac{1}{2}-1 \right )^2 + \left (0-1 \right )^2 = \dfrac{1}{4}+\dfrac{1}{4}+1 = \dfrac{6}{4}$ donc $\text{GJ}=\dfrac{\sqrt{6}}{2}$

[BD] est la diagonale du carré ABCD de côté 1 donc $\text{BD}=\sqrt{2}$.

L'aire de BDG est donc: $\mathcal A= \dfrac{\sqrt{2}\times \frac{\sqrt{6}}{2}}{2}=\dfrac{\sqrt{12}}{4}=\dfrac{\sqrt{3}}{2}$.
	\end{enumerate}

\item 
Le volume du tétraèdre EGBD est:
$\dfrac{1}{3} \left (\mathcal A \times \text{EI} \right )
= \dfrac{1}{3}\left ( \dfrac{\sqrt{3}}{2}\times \dfrac{2\sqrt{3}}{3}\right )
=\dfrac{1}{3}$


\end{enumerate}
